\documentclass{amsart}

\usepackage{amsmath}
\usepackage{amsthm}
\usepackage{amsfonts}
\usepackage{latexsym}
\usepackage{mathtools}
\usepackage{amssymb}
\usepackage{tikz}

\numberwithin{equation}{section}

\newcommand{\cat}{\text{cat}}
\newcommand{\ccat}{\text{ccat}}
\newcommand{\Int}{\mathrm{Int}}
\newcommand{\R}{\mathbb{R}}
\newcommand{\Z}{\mathbb{Z}}
\newcommand{\Nov}{\text{Nov}}
\newcommand{\grad}{\text{grad}}
\newcommand{\TC}{\text{TC}}
\newcommand{\GC}{\text{GC}}
\newcommand{\MH}{\mathrm{MH}}
\newcommand{\secat}{\text{secat}}

\newtheorem{theorem}{Theorem}[section]
\newtheorem{lemma}[theorem]{Lemma}
\newtheorem{definition}[theorem]{Definition}
\newtheorem{proposition}[theorem]{Proposition}
\newtheorem{corollary}[theorem]{Corollary}
\newtheorem{remark}[theorem]{Remark}
\newtheorem{example}[theorem]{Example}

\usepackage{docmute}
\usepackage[
  unicode,
  hidelinks,
  linktoc=all,
  bookmarksnumbered,
  hypertexnames=false
]{hyperref}

\hypersetup{
  pdftitle={
    Topological complexity, simplicial volume,
    and multiplicity of minimizing geodesics
  },
  pdfauthor={Kenji Fukushi}
}

\pdfstringdefDisableCommands{
  \def\({}
  \def\){}
  \def\Sigma{Sigma}
  \def\times{ x }
}

\title[Topological complexity and minimizing geodesics]{
  Topological complexity, simplicial volume,
  and multiplicity of minimizing geodesics
}

\author{Kenji Fukushi}

\subjclass[2020]{Primary 55M30; Secondary 53C22, 53C23, 20J06}

\keywords{Topological complexity, geodesic complexity, simplicial volume,
bounded cohomology, minimizing geodesics, magnitude homology}

\date{}

\begin{document}

\begin{abstract}
Let $M$ be a connected closed oriented $n$-manifold admitting a negatively
curved metric, with $n\geq2$. We prove that for an open dense
set of negatively curved metrics the number of ordered pairs
joined by exactly $2n+1$ minimizing geodesics is at least
$c_n\|M\|^2$, where $\|M\|$ is the simplicial volume and
$c_n>0$ is explicit.
The proof combines a localization formula for the exterior
canonical class of topological complexity with a quantitative
seminorm estimate.
For the latter we establish a norm-controlled lifting theorem
from group bounded cohomology to invariant cochains on a coset
space, assuming amenability of the pointwise stabilizers of
pairs of distinct cosets.
The local terms also admit an interpretation in minimal-length
magnitude homology.
Independently, a continuous selection argument gives
$\GC(M,g)=\TC(M)=2n+1$ in the unreduced convention for the
same open dense set of metrics.
\end{abstract}

\maketitle

\section{Introduction}

Topological complexity, introduced by Farber \cite{Farber2003}, measures the discontinuity of motion planning.
For a path-connected space $X$, it is the (unreduced) sectional
category of the endpoint fibration
\[
 \pi:PX\longrightarrow X\times X,\qquad
 \pi(\gamma)=(\gamma(0),\gamma(1)).
\]
For Riemannian manifolds, geodesic complexity imposes the additional
requirement that the chosen paths be minimizing geodesics
\cite{RecioMitter2021,MescherStegemeyer2023}.

Let $(M,g)$ be a closed negatively curved manifold of dimension
$n\geq 2$. The universal covering space $\widetilde M$ is a Hadamard
manifold, so fixed-endpoint homotopy classes determine unique geodesics.
Different classes can nevertheless contain minimizing geodesics of
equal length. Write $m_g(x,y)$ for their number.

We study how the multiplicity of minimizing geodesics reflects the
topology of $M$. Our primary goal is a quantitative lower bound,
involving simplicial volume, for the number of pairs admitting the
maximal number of minimizing geodesics under a general-position
condition. The topological complexity of closed negatively curved
manifolds is already known \cite{Dranishnikov2020}; our aim is not to reprove its value, but
to refine the geometric information associated with motion planning.

\subsection{Generic multiplicity}

Let $\mathcal R_-(M)$ denote the space of smooth negatively curved
metrics, with the $C^\infty$ topology, and put $\Gamma=\pi_1(M)$.
For lifts $u,v\in\widetilde M$ of $x,y\in M$ and $a\in\Gamma$, set
\[
 L_a(u,v)=d_{\widetilde M}(u,av),\qquad
 S_g(u,v)=\{a\in\Gamma:L_a(u,v)=d_M(x,y)\}.
\]
The set of minimizing labels $S_g(u,v)$ is finite, and
$m_g(x,y)=|S_g(u,v)|$.

\begin{definition}
A negatively curved metric $g$ is called \emph{regular} if, whenever
$S_g(u,v)=\{a_0,\ldots,a_{k-1}\}$ has $k\geq2$ elements, the map
\[
 (L_{a_1}-L_{a_0},\ldots,L_{a_{k-1}}-L_{a_0})
\]
has differential of rank $k-1$ at $(u,v)$.
\end{definition}

For a regular metric, $m_g(x,y)\leq2n+1$. Let
\[
 C_g=\{(x,y)\in M\times M:m_g(x,y)=2n+1\}.
\]
We regard these as ordered pairs.

\begin{theorem}\label{thm:intro-count}
Let $M$ be a connected, closed, oriented smooth $n$-manifold,
$n\geq2$, admitting a negatively curved metric.
Regular metrics form an open dense subset of $\mathcal R_-(M)$.
For each regular metric $g$, the set $C_g$ is finite and
\[
 \#C_g\geq c_n\|M\|^2,\qquad
 c_n=\begin{cases}
 \displaystyle\frac{2}{n^2((n+1)!)^2},&n\text{ even},\\[5pt]
 \displaystyle\frac{1}{2n(2n+1)!},&n\text{ odd}.
 \end{cases}
\]
Here $\|M\|$ denotes simplicial volume in the sense of
Gromov \cite{Gromov1982}.
\end{theorem}

\subsection{Exterior canonical classes and localization}

Put $G=\Gamma\times\Gamma$ and let $P=\Z[\Gamma]$ carry the action
$(g,h)\cdot a=gah^{-1}$. The augmentation kernel
$I=\ker(P\to\Z)$ supports the Costa--Farber canonical class
$\mathfrak v\in H^1(G;I)$ \cite{CostaFarber2010}. If $q_k:I^{\otimes k}\to\bigwedge^k I$
is the exterior quotient, define
\[
 \mathfrak a_k=q_{k*}(\mathfrak v^k),\qquad
 A_M=\mathfrak a_{2n}\frown[M\times M]
 \in(\bigwedge^{2n}I)_G.
\]
On $V_k=(\bigwedge^k(I\otimes\R))_G$, we use the quotient
$\ell^1$-seminorm $\|\cdot\|_\Delta$ whose generators of cost one are
\[
 [(a_1-a_0)\wedge\cdots\wedge(a_k-a_0)].
\]

\begin{theorem}\label{thm:intro-estimate}
For $M$ as in Theorem~\ref{thm:intro-count},
\[
 c_n\|M\|^2\leq\|A_M\|_\Delta\leq\|M\times M\|.
\]
\end{theorem}

The central geometric construction is the nerve of the sets of
simultaneously minimizing geodesic labels. For a regular metric its
top-dimensional simplices correspond to the pairs in $C_g$.
We obtain a localization formula expressing $A_M$ as a signed sum
of the associated exterior augmentation generators. Consequently
$\|A_M\|_\Delta\leq\#C_g$, which combines with
Theorem~\ref{thm:intro-estimate} to give
Theorem~\ref{thm:intro-count}.

These local terms have a metric-homological interpretation.
We use magnitude homology of metric spaces as developed in
\cite{LeinsterShulman2021}; its relationship with interval
order complexes is studied in \cite{KanetaYoshinaga2021}.
For distinct $x,y\in M$, the endpoint-fixed magnitude homology
at the minimal length satisfies
\[
 \MH_{2,d_M(x,y)}(M;x,y)\cong
 \ker(\Z[S_g(u,v)]\longrightarrow\Z).
\]
In particular, its rank equals $m_g(x,y)-1$; at pairs in $C_g$
the top exterior power is generated by the corresponding local
term in the localization formula.

\subsection{Bounded cohomology and geodesic complexity}

The estimate for $\|A_M\|_\Delta$ uses bounded cohomology and
its relation to simplicial volume \cite{Gromov1982}.
For the $G$-set $\Gamma\cong G/\Delta\Gamma$, the common stabilizer
of two distinct labels is a cyclic group: it is isomorphic to the
centralizer of a nontrivial element of $\Gamma$. This amenability
property permits controlled lifts of certain bounded classes
from $G$ to bounded invariant alternating cochains on $\Gamma$.

For regular metrics, smooth multiplicity strata also yield the
geodesic-complexity bound
\[
 \GC(M,g)\leq 2n+1.
\]
Together with the known value $\TC(M)=2n+1$
(in our unreduced convention) \cite{Dranishnikov2020}, this gives
\[
 \GC(M,g)=\TC(M)=2n+1
 =1+\max_{x\ne y}\operatorname{rank}
 \MH_{2,d_M(x,y)}(M;x,y).
\]
Here $\GC$ is defined using the unreduced convention.

The remainder of the paper is organized as follows. Section 2 studies canonical classes and magnitude homology.
Section 3 introduces the minimizing-geodesic nerve.
Section 4 develops the bounded cohomological lifting argument,
and Section 5 proves quantitative seminorm estimates.
Section 6 treats generic metrics and geodesic complexity.
Section 7 establishes localization and the counting theorem.
Section 8 discusses further questions.


\section{Canonical classes and magnitude homology}
\label{sec:canonical}

Throughout the paper $M$ is a connected closed negatively curved
Riemannian manifold of dimension $n\geq2$, $\Gamma=\pi_1(M)$, and
$G=\Gamma\times\Gamma$. We use integral coefficients unless indicated
otherwise. The definitions in this section require considerably less
geometry than negative curvature.

\subsection{The Costa--Farber class}

The endpoint fibration of an aspherical space can be compared with
the covering corresponding to the diagonal subgroup
$\Delta\Gamma\leq G$; see \cite{CostaFarber2010}.
Let $P=\Z[\Gamma]$ with the $G$-action
\[
 (g,h)a=gah^{-1}.
\]
Write $I=\ker(\varepsilon:P\to\Z)$.
The augmentation sequence is an extension of $G$-modules:
\begin{equation}\label{eq:augmentation}
 0\longrightarrow I\longrightarrow P
 \xrightarrow{\varepsilon}\Z\longrightarrow0.
\end{equation}
Its connecting class is
$\mathfrak v\in H^1(G;I)$. With the inhomogeneous convention
for group cochains, $\mathfrak v$ is represented by
\begin{equation}\label{eq:canonical-cocycle}
 c(g,h)=(g,h)\cdot1-1=gh^{-1}-1.
\end{equation}
In particular, the restriction of $\mathfrak v$ to
$\Delta\Gamma$ vanishes.

For each $k\geq1$, let
$q_k:I^{\otimes k}\to\bigwedge^k I$ be the quotient
mapping pure tensors to exterior products. Define
\begin{equation}\label{eq:exterior-class}
 \mathfrak a_k=q_{k*}(\mathfrak v^k)
 \in H^k(G;\bigwedge^k I).
\end{equation}
We emphasize that $q_k$ is the exterior quotient, not the
unnormalized alternating sum. No factorial occurs in
\eqref{eq:exterior-class}.

\subsection{The full simplex and exterior augmentation}

Let $Y$ be the full simplicial complex on the vertex set $\Gamma$.
Its oriented simplicial chain group in degree $k$ is
\[
 C_k(Y;\Z)=\bigwedge^{k+1}P.
\]
The augmentation of $P$ induces a contraction
\[
 \iota_\varepsilon:\bigwedge^{k+1}P
 \longrightarrow\bigwedge^kP,\quad
 \iota_\varepsilon(a_0\wedge\cdots\wedge a_k)
 =\sum_{i=0}^{k}(-1)^i
   a_0\wedge\cdots\widehat{a_i}\cdots\wedge a_k,
\]
which is exactly the simplicial boundary.

\begin{lemma}\label{lem:exterior-quotient}
For $k\geq1$ there is a natural $G$-equivariant exact sequence
\[
 C_{k+1}(Y)\xrightarrow{\partial}C_k(Y)
 \xrightarrow{T_k}\bigwedge^kI\longrightarrow0,
\]
where
\[
 T_k[a_0,\ldots,a_k]
 =(a_1-a_0)\wedge\cdots\wedge(a_k-a_0).
\]
Consequently
$C_k(Y)/\partial C_{k+1}(Y)\cong\bigwedge^kI$.
\end{lemma}

\begin{proof}
Choose any $a_*\in\Gamma$ and split the underlying abelian
group as $P=\Z a_*\oplus I$. Then
\[
 \bigwedge^{k+1}P
 = (a_*\wedge\bigwedge^kI)\oplus\bigwedge^{k+1}I.
\]
Contraction with $\varepsilon$ sends $a_*\wedge w$ to $w$
and vanishes on $\bigwedge^{k+1}I$.
The image of the contraction in degree $k+1$ is
$\bigwedge^{k+1}I$, which is precisely the kernel of
contraction in degree $k$. Thus contraction induces the
claimed quotient. The displayed expression for $T_k$
follows by expressing all vertices as
$a_0+(a_i-a_0)$ and expanding the wedge.
Although the splitting is not equivariant, the resulting map is.
\end{proof}

Write $Q=|Y|/G$ and
$\mathcal V_k=(\bigwedge^k I)_G$.
The action has no inversions on simplices (using the unique-root
property of torsion-free hyperbolic groups
\cite{BridsonHaefliger1999}): if an element
permuted a finite set of vertices, some positive power would
fix each of them; the unique-root property for torsion-free
hyperbolic groups then forces each vertex to have been fixed
already. The cellular chains of $Q$ therefore identify with
$C_k(Y)_G$. Passing to coinvariants in the lemma gives a
map from cellular cycles to $\mathcal V_k$.

\begin{lemma}\label{lem:cycle-injective}
The assignment
\[
 \iota_k:H_k(Q;\Z)\longrightarrow\mathcal V_k,\qquad
 [z]\longmapsto[T_k(z)]
\]
is well-defined and injective for $k\geq1$.
\end{lemma}

\begin{proof}
The definition is independent of boundaries by
$T_k\partial=0$. Exactness of coinvariants applied to
Lemma~\ref{lem:exterior-quotient} identifies
$\mathcal V_k$ with
$C_k(Y)_G/\partial C_{k+1}(Y)_G$.
The homology group $H_k(Q)$ is the subgroup of that
quotient represented by cycles, proving injectivity.
\end{proof}

\begin{proposition}\label{prop:comparison}
Let $B=M\times M$ and suppose $M$ is oriented.
If $f:B\to Q$ is the map obtained from a $G$-equivariant
map $\widetilde M^2\to Y$, then, for each $k$,
the image under $\iota_k$ of the homological comparison
$f_*:H_k(B)\to H_k(Q)$ is the cap-product evaluation
of $\mathfrak a_k$, with the usual local-coefficient
identifications. In particular,
\begin{equation}\label{eq:A-comparison}
 \iota_{2n}f_*[B]
 =\mathfrak a_{2n}\frown[B].
\end{equation}
\end{proposition}

\begin{proof}
The homogeneous bar construction gives an explicit
comparison map from a free $G$-resolution to the augmented
simplex chains: a bar simplex $(g_0,\ldots,g_k)$ is sent
to $[g_0\cdot1,\ldots,g_k\cdot1]$.
Its image under $T_k$ is
\[
 (g_1\cdot1-g_0\cdot1)\wedge\cdots\wedge
 (g_k\cdot1-g_0\cdot1).
\]
The cup product of the cocycle \eqref{eq:canonical-cocycle}
evaluates on the corresponding bar simplex as the wedge
of consecutive differences. Their transformation into
the differences from the first vertex is unit triangular,
so their wedge products coincide. This proves equality
on bar chains; any equivariant map from a free resolution
is equivariantly chain-homotopic to the comparison map.
The cap-product assertion follows.
\end{proof}

\subsection{Magnitude homology at the shortest length}

We use the magnitude chain complex of Leinster--Shulman
\cite{LeinsterShulman2021}. For endpoints $x,y$ and a
length $\ell$, write $C_{k,\ell}(M;x,y)$ for the free
abelian group generated by proper chains
\[
 (x=x_0,x_1,\ldots,x_k=y),
 \qquad \sum_{i=1}^k d(x_{i-1},x_i)=\ell.
\]
The differential deletes an internal point only when it
lies metrically between its neighbors. See also the
interval-poset description in \cite{KanetaYoshinaga2021}.

\begin{proposition}\label{prop:MH-minimal}
Let $x\ne y$ in a complete negatively curved manifold $M$,
and assume the number $m_g(x,y)$ of minimizing geodesics
is finite. Then there is an isomorphism
\[
 \MH_{2,d(x,y)}(M;x,y)
 \cong\ker\bigl(\Z[S_g(u,v)]
               \xrightarrow{\varepsilon}\Z\bigr).
\]
In particular,
\[
 \operatorname{rank}\MH_{2,d(x,y)}(M;x,y)
 =m_g(x,y)-1.
\]
\end{proposition}

\begin{proof}
Write $d=d(x,y)>0$. A chain of total length $d$
must be ordered along a minimizing geodesic, by equality
in the triangle inequality. A point $z\ne x,y$ occurring
in $(x,z,y)$ determines exactly one minimizing geodesic:
the two minimizing segments splice without a corner,
because their concatenation is globally minimizing,
and the geodesic equation has unique solutions.
Conversely every minimizing geodesic contains such points.

Now $C_{1,d}(M;x,y)=\Z(x,y)$, and
\[
 \partial_2(x,z,y)=-(x,y).
\]
The differential in degree three on a chain
$(x,z,w,y)$, where $z$ precedes $w$ along a
minimizing geodesic, is
\[
 \partial_3(x,z,w,y)
 =-(x,w,y)+(x,z,y).
\]
Every pair of internal points on the same minimizing
geodesic can be related using such boundaries, taking
a common later point if necessary. There are no
degree-three boundaries relating interior points of
distinct minimizing geodesics, because all chains
of total length $d$ lie on a single minimizing
geodesic. Hence
\[
 C_{2,d}/\operatorname{im}\partial_3
 \cong\Z[S_g(u,v)]
\]
and the induced differential is minus augmentation.
Taking its kernel proves the result.
\end{proof}

For a minimizing-label set $S$ write
$I(S)=\ker(\Z[S]\to\Z)$. The proposition identifies
the local coefficient $I(S)$ with magnitude homology;
the inclusion $I(S)\hookrightarrow I$ is natural
after lifts have been chosen. This is the local-to-global
map used in the localization theorem below.

\begin{remark}
The endpoint-fixed statement is made only for $x\ne y$.
The diagonal has length zero and a different chain complex.
No identification of magnitude homology with the entire
simplicial nerve is asserted: the correspondence here
is specifically in degree two and at the minimal length.
\end{remark}


\section{The minimizing-geodesic nerve}
\label{sec:nerve}

Let $X=\widetilde M$ with its lifted negatively curved metric,
and put $B=M\times M$, $\widetilde B=X\times X$.
The group $G=\Gamma\times\Gamma$ acts on $\widetilde B$ by
\[
 (g,h)\cdot(u,v)=(gu,hv)
\]
and on the label set $\Gamma$ by
$(g,h)\cdot a=gah^{-1}$.
For $a\in\Gamma$, let
\[
 L_a(u,v)=d_X(u,av),\qquad
 m(u,v)=\min_{a\in\Gamma}L_a(u,v)=d_M(x,y).
\]
All minima are attained, and the set of minimizing labels
$S(u,v)=\{a:L_a(u,v)=m(u,v)\}$ is finite.

\subsection{The nerve and a uniform allowance}

For every label $a$, define the closed subset
\[
 D_a=\{(u,v)\in\widetilde B:L_a(u,v)=m(u,v)\}.
\]
The \emph{minimizing-geodesic nerve} $\mathcal N_g$
has vertex set $\Gamma$ and a simplex on each finite
nonempty $S\subset\Gamma$ with
$\bigcap_{a\in S}D_a\ne\varnothing$.
Equivariance of $L_a$ gives a simplicial $G$-action.

For $\delta>0$, also define the open sets
\[
 D_a^\delta=\{(u,v):L_a(u,v)<m(u,v)+\delta\}
\]
and let $\mathcal N_g^\delta$ be their nerve.
Clearly $\mathcal N_g\subseteq\mathcal N_g^\delta$.

\begin{lemma}\label{lem:uniform-allowance}
There exists $\delta_0>0$ such that
\[
 \mathcal N_g^\delta=\mathcal N_g
 \quad\text{for every }0<\delta\leq\delta_0.
\]
Moreover, $\mathcal N_g$ has finitely many $G$-orbits of simplices.
\end{lemma}

\begin{proof}
Choose a compact fundamental region $K\subset X$.
For $(u,v)\in K\times K$ and $L_a(u,v)\leq m(u,v)+1$,
the label $a$ belongs to the finite set
\[
 F=\{a\in\Gamma:d_X(K,aK)\leq\operatorname{diam}(M)+1\}.
\]
Thus every simplex of $\mathcal N_g^\delta$,
$0<\delta\leq1$, is $G$-equivalent to a subset
of the fixed finite set $F$.
For a fixed nonempty subset $S\subset F$, put
\[
 \theta_S(u,v)=\max_{a\in S}\bigl(L_a(u,v)-m(u,v)\bigr).
\]
If $S$ does not span a simplex of $\mathcal N_g$,
the continuous function $\theta_S$ is strictly positive
on the compact set $K\times K$, hence has a positive minimum.
There are only finitely many such $S$. Taking $\delta_0$
below all their positive minima proves the equality
of nerves. The same finite-label argument yields the
finiteness of simplex orbits.
\end{proof}

\begin{remark}
The inequality in the definition of $D_a^\delta$ is strict.
Therefore the value of $\delta_0$ can be reduced so that
the proof works without a boundary equality issue.
The lemma concerns the nerves of the covering sets,
not equality of the sets $D_a$ and $D_a^\delta$ themselves.
\end{remark}

\subsection{The equivariant map}

Choose $0<\delta\leq\delta_0$ and define
\[
 w_a(u,v)=\max\{0,m(u,v)+\delta-L_a(u,v)\},
 \qquad W(u,v)=\sum_{a\in\Gamma}w_a(u,v).
\]
At every $(u,v)$ at least one label minimizes $L_a$,
so $W(u,v)>0$.

\begin{proposition}\label{prop:Phi}
The expression
\begin{equation}\label{eq:Phi}
 \Phi_g(u,v)=\sum_{a\in\Gamma}
 \frac{w_a(u,v)}{W(u,v)}[a]
\end{equation}
defines a continuous $G$-equivariant map
\[
 \Phi_g:\widetilde B\longrightarrow|\mathcal N_g|.
\]
It descends to a continuous map
\[
 \bar\Phi_g:B\longrightarrow Q_g=|\mathcal N_g|/G.
\]
\end{proposition}

\begin{proof}
On any compact subset of $\widetilde B$, the label set
for which $w_a$ can be nonzero is finite, by proper
discontinuity of the deck action. Thus \eqref{eq:Phi}
is a locally finite sum of continuous functions.
Its support at $(u,v)$ is a simplex of
$\mathcal N_g^\delta=\mathcal N_g$, hence its image
belongs to the geometric realization.
The identity
\[
 L_{gah^{-1}}(gu,hv)=L_a(u,v)
\]
and invariance of $m$ establish equivariance.
The quotient statement follows.
\end{proof}

\begin{lemma}\label{lem:no-inversions}
The $G$-action on $\mathcal N_g$ is without inversions:
every element preserving a simplex setwise fixes its vertices.
Consequently $Q_g$ has a finite CW structure with one
open cell for each $G$-orbit of simplices.
\end{lemma}

\begin{proof}
Suppose $(g,h)$ preserves a finite vertex set $S$.
For every $a\in S$, some common positive power $r$
of $(g,h)$ fixes $a$, so $g^r=ah^ra^{-1}$.
The torsion-free hyperbolic group $\Gamma$ has the
unique-root property: $p^r=q^r$ implies $p=q$.
It follows that $g=aha^{-1}$ for every $a\in S$,
so $(g,h)$ fixes every vertex.
The standard quotient-cell construction gives the
last assertion, using Lemma~\ref{lem:uniform-allowance}.
\end{proof}

The quotient $Q_g$ need not be contractible. Its relevance
comes instead from the fundamental cycle represented by
$\bar\Phi_g$ and its image in the exterior augmentation
module under Lemma~\ref{lem:cycle-injective}.

\subsection{Multiplicity and the dimension of the nerve}

If $g$ is regular in the sense of Section 1, every minimizing
label set has cardinality at most $2n+1$. Consequently,
\[
 \dim\mathcal N_g\leq2n.
\]
At a point with precisely $k$ minimizing labels, the
length-difference equations cut out a smooth local stratum
of codimension $k-1$ in $B$, by the implicit function
theorem. The inequalities excluding all remaining labels
are strict in a sufficiently small neighborhood.
Thus the local geometry of each simplex is controlled
by the distance-difference functions, even though the
global nerve is defined only in terms of closed sets.

\begin{proposition}\label{prop:top-simplex}
Suppose $g$ is regular and set $N=2n$.
Every pair in $C_g$ determines an $N$-simplex
of $\mathcal N_g$. Conversely, every $N$-simplex
is realized by one or more isolated endpoint pairs.
There are only finitely many such pairs modulo $G$,
and thus $C_g$ is finite.
\end{proposition}

\begin{proof}
An $N$-simplex has $N+1$ minimizing labels.
Its defining $N$ length-difference equations have
full-rank differential at each realization, so their
solutions are isolated. To see that the set of
realizations modulo $G$ is finite, put endpoints
in compact fundamental regions. By properness,
only finitely many labels can minimize there.
If infinitely many pairs realized top simplices,
a convergent subsequence would have a fixed
label configuration and converge to another
simultaneous minimum. Full-rank transversality
at that limit contradicts accumulation.
Each pair in $C_g$ has exactly $N+1$ labels,
so it determines exactly one top simplex after lifts.
Different pairs can determine the same simplex
orbit; no global bijection with simplex orbits
is asserted.
\end{proof}

\begin{remark}\label{rem:nerve-geometry}
The full simplex $Y$ from Section~\ref{sec:canonical}
is contractible and contains $\mathcal N_g$ equivariantly.
The composite
$B\xrightarrow{\bar\Phi_g}Q_g\to Q$
induces, on the fundamental class, the element
$A_M$ of \eqref{eq:A-comparison}. We do not require
the closed sets $D_a$ to constitute a good cover.
\end{remark}


\section{A bounded cohomological lifting theorem}
\label{sec:bounded}

The quantitative estimate will use bounded cochains on a
transitive $G$-set. We give a norm-controlled lifting argument.
The construction is related to resolutions and relative bounded
cohomology; compare the bounded-cohomological framework of
\cite{Gromov1982}. The treatment of equivariant cochain complexes and boundedly acyclic
actions in \cite[Corollary~5.15 and Remark~5.17]{LiLohMoraschini2026}
is closely related to our construction. Unlike a globally amenable
action, the present coset action has a nonamenable vertex stabilizer;
our argument retains this single-column contribution and gives an
explicit cochain-level bound. We make no claim that the general
lifting mechanism is new.

We introduce the homogeneous and coset cochain complexes.

Let $G$ be a discrete group, $H\leq G$, and $S=G/H$.
The homogeneous complex of bounded real cochains on $G$ is
\[
 C_b^q(G)=\ell^\infty(G^{q+1};\R)^G.
\]
We also consider
\[
 A^p=\ell^\infty_{\mathrm{alt}}(S^{p+1};\R)^G
\]
with the usual alternating simplicial coboundary $\delta$.
Pullback along $\pi:G\to S$ induces a cochain map
$A^\bullet\to C_b^\bullet(G)$, hence
\[
 \rho^k:H^k(A^\bullet)\longrightarrow H_b^k(G;\R).
\]

For all $p,q\geq0$, define
\[
 D^{p,q}=
 \ell^\infty_{\mathrm{alt},S}
 (S^{p+1}\times G^{q+1};\R)^G,
\]
where alternation is imposed only on the $S$-variables.
The horizontal and vertical coboundaries are $\delta$ and $d_G$,
respectively; the total differential on $D^{p,q}$ is
$\delta+(-1)^p d_G$.

\begin{lemma}\label{lem:AW-comparison}
For $p,q\geq0$, let
\[
 R_{p,q}:D^{p,q}\longrightarrow C_b^{p+q}(G)
\]
be the bounded linear map defined by
\[
 (R_{p,q}f)(g_0,\ldots,g_{p+q})
 =f(g_0H,\ldots,g_pH;g_p,\ldots,g_{p+q}).
\]
The maps $R_{p,q}$ have operator norm at most one,
and their sum $R$ on the total complex satisfies
\[
 d_G R=R\bigl(\delta+(-1)^p d_G\bigr).
\]
For $f\in D^{k,0}$ satisfying $d_Gf=0$,
$R f$ is precisely the coset pullback of the
$G$-invariant alternating function determined by $f$.
\end{lemma}

\begin{proof}
Boundedness and $G$-invariance are immediate.
To check the cochain identity, evaluate both sides on
$(g_0,\ldots,g_{p+q+1})$. The horizontal differential
contributes
\[
 \sum_{i=0}^{p+1}(-1)^i
 f(g_0H,\ldots,\widehat{g_iH},\ldots,g_{p+1}H;
   g_{p+1},\ldots,g_{p+q+1}),
\]
and the signed vertical differential contributes
\[
 (-1)^p\sum_{j=0}^{q+1}(-1)^j
 f(g_0H,\ldots,g_pH;
   g_p,\ldots,\widehat{g_{p+j}},\ldots,g_{p+q+1}).
\]
The horizontal term $i=p+1$ cancels the vertical
term $j=0$. All remaining terms are exactly
the homogeneous coboundary of $R_{p,q}f$:
the first $p+1$ delete vertices in the
$S$-variables and the rest delete group variables.
This proves the identity.

If $d_Gf=0$ in vertical degree zero,
$f(s_0,\ldots,s_k;g)$ is independent of $g$,
so $R_{k,0}f$ is the pullback of the corresponding
invariant alternating function of the coset vertices.
\end{proof}

\begin{lemma}\label{lem:horizontal-contraction}
For every $q\geq0$ the horizontally augmented
complex
\[
 0\longrightarrow C_b^q(G)
 \xrightarrow{\jmath}D^{0,q}
 \xrightarrow{\delta}D^{1,q}
 \xrightarrow{\delta}D^{2,q}\longrightarrow\cdots
\]
is contractible by maps of operator norm at most one.
Consequently the augmentation
\[
 \jmath:C_b^\bullet(G)\longrightarrow
          \operatorname{Tot}(D^{\bullet,\bullet}),
 \qquad (\jmath\alpha)(s;g_0,\ldots,g_q)
          =\alpha(g_0,\ldots,g_q)
\]
is a quasi-isomorphism. Under this identification,
the cochain map $R$ of Lemma~\ref{lem:AW-comparison}
is its inverse on cohomology.
\end{lemma}

\begin{proof}
For $p\geq1$ define
\[
 (h^{p,q}f)(s_0,\ldots,s_{p-1};
                         g_0,\ldots,g_q)
 =
 f(g_0H,s_0,\ldots,s_{p-1};g_0,\ldots,g_q).
\]
For $p=0$ use the same formula with no $s$-entries:
\[
 (h^{0,q}f)(g_0,\ldots,g_q)
 =f(g_0H;g_0,\ldots,g_q).
\]
The maps are $G$-equivariant and preserve alternation
in the remaining $S$-variables; their norms are at
most one. Direct expansion of the alternating
coboundary gives
\[
 h^{p+1,q}\delta+\delta h^{p,q}
 =\mathrm{id}_{D^{p,q}}\quad(p\geq1),
 \qquad
 h^{1,q}\delta+\jmath h^{0,q}
 =\mathrm{id}_{D^{0,q}},
 \qquad h^{0,q}\jmath=\mathrm{id}.
\]
The formulas involve $g_0H$, which is why
the homotopies are equivariant despite the
absence of a $G$-fixed vertex of $S$.

Thus the horizontal augmentation is contractible
in every fixed vertical degree. Filtering the
first-quadrant total complex by the vertical
degree, its horizontal-first spectral sequence
has only the augmented-column contribution
$C_b^\bullet(G)$; equivalently, the usual
bounded double-complex comparison lemma
applies to this explicitly contractible
horizontal resolution. This proves that $\jmath$
is a quasi-isomorphism. Finally
$R\jmath=\mathrm{id}_{C_b^\bullet(G)}$
on cochains, so $R$ induces the inverse.
\end{proof}

\begin{theorem}\label{thm:bounded-lift}
Assume that the pointwise stabilizer of any two distinct
elements of $S$ is amenable. Then for each $k\geq2$
\begin{equation}\label{eq:bounded-image}
 \operatorname{im}\rho^k
 =\ker\bigl(H_b^k(G;\R)\longrightarrow H_b^k(H;\R)\bigr).
\end{equation}
More quantitatively, if $\alpha\in C_b^k(G)$ is a
homogeneous cocycle satisfying
$\alpha|_{H^{k+1}}=0$, there is a cocycle $\beta\in A^k$
such that
\[
 [\pi^*\beta]=[\alpha],\qquad
 \|\beta\|_\infty\leq k(k+1)!\|\alpha\|_\infty.
\]
\end{theorem}

We first record the bounded contractions used in the proof.

\begin{lemma}\label{lem:bounded-vertical}
For $p\geq1$, the vertical homogeneous coboundary complex
$D^{p,\bullet}$ has bounded contracting homotopies
in positive degrees, of norm at most one.
\end{lemma}

\begin{proof}
Consider an ordered tuple of pairwise distinct vertices
$(s_0,\ldots,s_p)$. Its pointwise stabilizer $K$ fixes
at least two distinct vertices, so is amenable.
The setwise stabilizer is a finite extension of $K$
and therefore also amenable. On each orbit of such
tuples, equivariant cochains are identified with bounded
cochains on $G^{q+1}$ invariant under the tuple stabilizer,
with the sign representation accounting for odd vertex
permutations.

We make the contraction explicit. Let $K_\sigma$
be the setwise stabilizer of the unordered vertex set
of an ordered tuple $\sigma$, and let
$\chi_\sigma:K_\sigma\to\{+1,-1\}$ be its permutation-sign
character. Fix a left $K_\sigma$-invariant mean
$m_\sigma$ on $\ell^\infty(G)$. Such a mean exists:
$K_\sigma$ is amenable, and one may choose a mean on
each $K_\sigma$-orbit of $G$ and then a mean on the
set of these orbits. On the vertical complex associated
with $\sigma$, cochains satisfy
\[
 f(kg_0,\ldots,kg_q)
  =\chi_\sigma(k)f(g_0,\ldots,g_q).
\]
Define, for $q\geq1$,
\[
 (h_\sigma f)(g_0,\ldots,g_{q-1})
  =m_\sigma\bigl(x\mapsto
       f(x,g_0,\ldots,g_{q-1})\bigr).
\]
Left invariance of $m_\sigma$ shows that
$h_\sigma f$ obeys the same $\chi_\sigma$-equivariance
as $f$. The standard insertion identity
\[
 d_Gh_\sigma+h_\sigma d_G=\mathrm{id}
\]
holds in positive vertical degrees, and
$\|h_\sigma\|\leq1$. The sign character is
handled by equivariance; no additional
averaging factor is required.

Choose an orbit representative for each unordered
tuple of distinct vertices and transport this
construction equivariantly. The resulting
homotopies have uniform norm at most one.
Alternation forces cochains on tuples with
repeated vertices to vanish.
\end{proof}

\begin{lemma}\label{lem:relative-primitive}
Let $\alpha\in C_b^k(G)$, $k\geq1$, be a bounded
homogeneous cocycle with zero restriction as a cochain
to $H^{k+1}$. The corresponding element of
$D^{0,k}$ admits a vertical primitive $b_0\in
D^{0,k-1}$ satisfying
\[
 d_Gb_0=\alpha,\qquad
 \|b_0\|_\infty\leq k\|\alpha\|_\infty.
\]
\end{lemma}

\begin{proof}
Choose an $H$-equivariant retraction of sets
$r:G\to H$ that is the identity on $H$; choose
representatives of the left $H$-orbits, with the
representative $1$ for $H$ itself. The standard
prism homotopy between the identity and $r$ on
the homogeneous bar complex has, on cochains
of degree $k$, at most $k$ summands. Consequently
its operator norm is at most $k$.
Since $\alpha$ restricts to zero on $H$,
the prism identity $dK+Kd=\mathrm{id}-r^*$
gives $d_G(K\alpha)=\alpha$.
Transporting this primitive from the vertex $H$
to the other vertices of $G/H$ gives the
required $b_0$.
\end{proof}

\begin{proof}[Proof of Theorem~\ref{thm:bounded-lift}]
We use the first-quadrant double complex
$D^{p,q}$ with horizontal differential
$\delta$ and vertical differential $d_G$, and
total differential $d_{\mathrm{tot}}=
\delta+(-1)^pd_G$ on $D^{p,q}$.
The horizontal contraction in
Lemma~\ref{lem:horizontal-contraction}
shows that $\jmath:C_b^\bullet(G)\to
\operatorname{Tot}D$ is a quasi-isomorphism.
The map $R$ of Lemma~\ref{lem:AW-comparison}
is a cochain left inverse to $\jmath$.
In particular, a cocycle in
\[
 A^p=\ker(d_G:D^{p,0}\to D^{p,1})
\]
has image under $R$ equal to its coset
pullback.

First suppose that $\alpha\in C_b^k(G)$
is a cocycle whose restriction to $H^{k+1}$
vanishes as a cochain. Write
$a_0=\jmath\alpha\in D^{0,k}$.
Lemma~\ref{lem:relative-primitive}
provides $b_0\in D^{0,k-1}$ with
\[
 d_Gb_0=a_0,\qquad
 \|b_0\|_\infty\leq k\|\alpha\|_\infty.
\]
Set $a_1=-\delta b_0\in D^{1,k-1}$.
Then $a_0-d_{\mathrm{tot}}b_0=a_1$,
and $\delta a_1=d_Ga_1=0$.

For $1\leq p<k$, suppose inductively
that $a_p\in D^{p,k-p}$ is closed under
both $\delta$ and $d_G$. Its vertical
degree $k-p$ is positive, so
Lemma~\ref{lem:bounded-vertical}
gives a primitive
\[
 b_p=(-1)^ph_pa_p\in D^{p,k-p-1},
 \qquad d_Gb_p=(-1)^pa_p,
 \qquad \|b_p\|_\infty\leq\|a_p\|_\infty.
\]
Put $a_{p+1}=-\delta b_p$. Then
\[
 a_p-d_{\mathrm{tot}}b_p=a_{p+1},
 \qquad
 \delta a_{p+1}=0,
 \qquad
 d_Ga_{p+1}=0.
\]
The last identity follows from
$d_G\delta=\delta d_G$ and
$\delta a_p=0$. By induction
$a_k\in D^{k,0}$ is closed under
both differentials, so it is an element
$\beta\in A^k$ with $\delta\beta=0$.
Summing the displayed total-coboundary
identities and applying $R$ proves
$[\pi^*\beta]=[\alpha]$.

On $D^{p,q}$ the horizontal coboundary
has operator norm at most $p+2$.
Consequently
\[
 \|\beta\|_\infty
 \leq (k+1)!\|b_0\|_\infty
 \leq k(k+1)!\|\alpha\|_\infty.
\]

Now suppose only that
$[\alpha]|_H=0$ in bounded cohomology.
There is a bounded homogeneous
$H$-invariant $(k-1)$-cochain
$\theta$ on $H$ with
$d_G\theta=\alpha|_{H^{k+1}}$.
Define a bounded homogeneous
$G$-invariant cochain $\widehat\theta$
on $G^k$ as follows. If all its
arguments lie in a common left coset
$tH$, set
\[
 \widehat\theta(th_0,\ldots,th_{k-1})
    =\theta(h_0,\ldots,h_{k-1});
\]
otherwise set its value to zero.
This is well-defined: replacing $t$
by $th$ translates all arguments
of $\theta$ simultaneously by $h^{-1}$,
and $\theta$ is $H$-invariant.
The resulting cochain is bounded
and $G$-invariant. Thus
$\alpha'=\alpha-d_G\widehat\theta$
has zero cochain restriction to $H$.
Applying the first part to $\alpha'$
gives a cocycle in $A^k$ mapping to
$[\alpha]$.

Conversely, any alternating cocycle on
$S$ pulls back to a cocycle vanishing
on $H^{k+1}$, because all those
arguments have coset $H$. Hence
$\operatorname{im}\rho^k$ is exactly
the restriction kernel, proving
\eqref{eq:bounded-image}.
\end{proof}

We specialize the lifting theorem to the diagonal subgroup.

For the application let $\Gamma$ be the fundamental group
of a closed negatively curved manifold and put
\[
 G=\Gamma\times\Gamma,\qquad H=\Delta\Gamma.
\]
The $G$-set $G/H$ is naturally isomorphic to $\Gamma$,
with $(g,h)a=gah^{-1}$.

\begin{lemma}\label{lem:diagonal-stabilizers}
The common stabilizer in $G$ of two distinct labels
$a,b\in\Gamma$ is isomorphic to
$C_\Gamma(b^{-1}a)$. Hence it is infinite cyclic.
\end{lemma}

\begin{proof}
An element $(g,h)$ fixes $a$ if and only if
$g=aha^{-1}$; it fixes $b$ as well exactly when
$h$ commutes with $b^{-1}a$.
Since $b^{-1}a\ne1$, its centralizer in the
torsion-free hyperbolic group $\Gamma$ is
infinite cyclic \cite{BridsonHaefliger1999}. This identifies the
common stabilizer.
\end{proof}

\begin{proposition}\label{prop:difference-lift-bound}
Let $\omega\in C_b^n(\Gamma)$ be an alternating homogeneous
$n$-cocycle, $n\geq2$, and put
\[
 D=p_2^*\omega-p_1^*\omega\in C_b^n(G).
\]
There is a bounded alternating $G$-invariant cocycle
$\widetilde D$ on $\Gamma^{n+1}$ with
\[
 [\pi^*\widetilde D]=[D],\qquad
 \|\widetilde D\|_\infty
 \leq n(n+1)!\,\|\omega\|_\infty.
\]
\end{proposition}

\begin{proof}
For the initial vertical contraction at the coset
$H=\Delta\Gamma$, use the $H$-equivariant retraction
\[
 r:G\longrightarrow H,\qquad r(g,h)=(g,g).
\]
The homogeneous prism cochain of degree $n-1$ is
\[
 KD(g_0,\ldots,g_{n-1})
 =\sum_{i=0}^{n-1}(-1)^i
 D(rg_0,\ldots,rg_i,g_i,\ldots,g_{n-1}).
\]
Here the expression in the $i$-th summand has
$n+1$ group arguments. The first-coordinate
parts of $rg_i$ and $g_i$ agree. Since $\omega$
is alternating, the $p_1^*\omega$ contribution
of every summand is zero. Consequently
\[
 \|KD\|_\infty\leq n\|\omega\|_\infty.
\]
The prism identity and $D|_{H^{n+1}}=0$
give $d_G(KD)=D$.

Transport $KD$ equivariantly to the other vertex
cosets, obtaining $b_0\in D^{0,n-1}$
with the same bound. In the double complex
of Theorem~\ref{thm:bounded-lift}, set
\[
 a_1=-\delta b_0,\qquad
 b_p=(-1)^p h_pa_p,\qquad
 a_{p+1}=-\delta b_p
 \quad(1\leq p<n),
\]
where $h_p$ denotes the norm-one vertical
contraction in column $p$. This removes the
positive vertical degrees successively.
The horizontal differential from column
$p$ to column $p+1$ has norm at most $p+2$.
It follows that the final cocycle
$\widetilde D=a_n$ satisfies
\[
 \|\widetilde D\|_\infty
 \leq n(2\cdot3\cdots(n+1))
       \|\omega\|_\infty
 =n(n+1)!\|\omega\|_\infty.
\]
The total-complex coboundary identities imply
$[\pi^*\widetilde D]=[D]$, as in
Theorem~\ref{thm:bounded-lift}.
\end{proof}

\begin{corollary}\label{cor:bounded-diagonal}
For $k\geq2$, every bounded cohomology class on
$\Gamma\times\Gamma$ restricting trivially to
$\Delta\Gamma$ is represented by the pullback
of a bounded alternating invariant cocycle
on $\Gamma^{k+1}$. A cocycle whose restriction
vanishes literally admits the norm estimate
of Theorem~\ref{thm:bounded-lift}.
\end{corollary}

\begin{remark}
The norm estimate concerns a cocycle with zero
cochain restriction. It is not a uniform bound
in terms of the norm of an arbitrary class in
the kernel unless a norm-controlled choice of
relative representative is supplied.
This distinction is important in Section 5,
where the cocycles are chosen to vanish
explicitly on the diagonal.
\end{remark}


\section{A quantitative exterior seminorm}
\label{sec:seminorm}

We prove the estimate in Theorem~\ref{thm:intro-estimate}.
Throughout this section $M$ is closed, oriented and negatively
curved, of dimension $n\geq2$. Write $B=M\times M$,
$G=\Gamma\times\Gamma$, $N=2n$, and
\[
 V_k=(\bigwedge^k I_{\R})_G,\qquad I_{\R}=I\otimes\R.
\]

We first compare the exterior seminorm with a cellular quotient seminorm.

For an ordered tuple of labels put
\[
 u(a_0,\ldots,a_k)
 =[(a_1-a_0)\wedge\cdots\wedge(a_k-a_0)]
 \in V_k.
\]
Define
\begin{equation}\label{eq:delta-norm}
 \|z\|_\Delta=
 \inf\left\{\sum_{j=1}^{r}|t_j|:
 z=\sum_{j=1}^{r}t_j
    u(a_{j,0},\ldots,a_{j,k})\right\}.
\end{equation}
The infimum is over finite expressions. This is a
possibly degenerate seminorm, not an asserted norm.

Let $Q=|Y|/G$, where $Y$ is the full simplex on
$\Gamma$. Endow the cellular chain complex of
$Q$ with the quotient $\ell^1$-seminorm assigning
cost one to each oriented simplex orbit.

\begin{lemma}\label{lem:isometric-image}
The injection
$\iota_k:H_k(Q;\R)\hookrightarrow V_k$
from Lemma~\ref{lem:cycle-injective} preserves
the quotient $\ell^1$-seminorm:
\[
 \|\iota_k[z]\|_\Delta
 =\inf\{\|c\|_1:\partial c=0,\ [c]=[z]\}.
\]
In particular, if $\beta:\Gamma^{k+1}\to\R$
is a bounded alternating $G$-invariant cocycle,
it induces a linear functional on $V_k$, and
\begin{equation}\label{eq:dual-pairing}
 |\langle\beta,z\rangle|
 \leq\|\beta\|_\infty\|z\|_\Delta.
\end{equation}
\end{lemma}

\begin{proof}
By Lemma~\ref{lem:exterior-quotient},
$V_k$ is the quotient of $C_k(Y;\R)_G$
by $\partial C_{k+1}(Y;\R)_G$. The generators
in \eqref{eq:delta-norm} are exactly the
images of the oriented simplex generators.
If $z$ is a cycle, every chain with the same
image in this quotient differs from it by a
boundary, and therefore is again a cycle.
Thus the infima defining both seminorms agree.

An alternating invariant cocycle vanishes
on cellular boundaries, so it defines a
functional on the quotient, of norm bounded
by its supremum norm. This proves
\eqref{eq:dual-pairing}.
\end{proof}

\begin{lemma}\label{lem:upper-norm}
For the class
$A_M=\mathfrak a_N\frown[B]$ one has
\[
 \|A_M\|_\Delta\leq\|B\|.
\]
\end{lemma}

\begin{proof}
Use a singular fundamental cycle of $B$ whose
$\ell^1$-norm is arbitrarily close to $\|B\|$.
The bar comparison followed by $T_N$ in
Proposition~\ref{prop:comparison} sends each
oriented singular simplex to an exterior
generator with coefficient of absolute
value at most one. Therefore the induced
chain map is nonexpanding in the $\ell^1$-norm.
Taking the infimum proves the inequality.
\end{proof}

The lower estimate is obtained from a bounded orientation cocycle.

By the bounded-cohomological dual description of
simplicial volume \cite{Gromov1982}, for every
$\epsilon>0$ there is a bounded alternating
homogeneous $n$-cocycle $\omega$ on $\Gamma$
representing the real orientation class,
normalized by
\begin{equation}\label{eq:orientation-norm}
 \langle[\omega],[M]\rangle=1,\qquad
 \|\omega\|_\infty\leq\|M\|^{-1}+\epsilon.
\end{equation}
Negative curvature ensures $\|M\|>0$.
Pull $\omega$ back to $G$ along the two
projections and denote the resulting cocycles by
$A=p_1^*\omega$ and $B'=p_2^*\omega$.
We write $B'$ to distinguish this cocycle
from the manifold $B=M\times M$.

Assume first that $n$ is odd.
Consider the $N$-cocycle on $G$
\[
 \alpha=\frac12(A\smile B'-B'\smile A).
\]
Its restriction to $\Delta\Gamma$ is identically
zero as a cochain. The cup-product sup-norm
estimate gives
\[
 \|\alpha\|_\infty\leq\|\omega\|_\infty^2.
\]
Since $n$ is odd, graded commutativity
in ordinary cohomology gives
\[
 \langle[\alpha],[M\times M]\rangle=1.
\]
By Theorem~\ref{thm:bounded-lift}, $\alpha$
lifts to a bounded alternating invariant
$N$-cocycle $\beta$ on $\Gamma^{N+1}$, with
\[
 \|\beta\|_\infty
 \leq N(N+1)!\,\|\omega\|_\infty^2.
\]
The bar comparison of Proposition
\ref{prop:comparison}, followed by
\eqref{eq:dual-pairing}, gives
\[
 1=|\langle\beta,A_M\rangle|
 \leq N(N+1)!\|\omega\|_\infty^2
       \|A_M\|_\Delta.
\]
Letting $\epsilon\downarrow0$ yields
\begin{equation}\label{eq:odd-bound}
 \|A_M\|_\Delta\geq
 \frac{\|M\|^2}{2n(2n+1)!}.
\end{equation}

We turn to the case where $n$ is even and set
\[
 D=B'-A\in C_b^n(G).
\]
As cochains, $D|_{\Delta\Gamma}=0$.
Theorem~\ref{thm:bounded-lift} supplies a
bounded alternating invariant cocycle
$\widetilde D$ on $\Gamma^{n+1}$
with $\pi^*[\widetilde D]=[D]$.

For the explicit bound we use the retraction
$r:G\to\Delta\Gamma$, $(g,h)\mapsto(g,g)$,
in the initial prism of Lemma
\ref{lem:relative-primitive}. In that prism,
the summands involving $p_1^*\omega$ vanish:
at the transition between $r(g_i,h_i)$ and
$(g_i,h_i)$ the first-coordinate arguments
repeat, and $\omega$ is alternating.
Hence only the $p_2^*\omega$ terms contribute.
The primitive therefore has norm at most
$n\|\omega\|_\infty$, rather than
$2n\|\omega\|_\infty$.
The remaining horizontal and vertical
contractions give
\begin{equation}\label{eq:Dtilde-bound}
 \|\widetilde D\|_\infty
 \leq n(n+1)!\|\omega\|_\infty.
\end{equation}

Define a bounded alternating invariant
$N$-cocycle on $\Gamma$ by
\[
 \beta=-\frac12
 \operatorname{Alt}(\widetilde D\smile\widetilde D),
\]
where $\operatorname{Alt}$ denotes normalized
alternation (average over permutations with
sign). Since normalized alternation has norm
at most one and induces the identity in
cohomology with real coefficients,
\[
 \|\beta\|_\infty
 \leq\frac12 n^2((n+1)!)^2
       \|\omega\|_\infty^2.
\]
The cohomology class of $\pi^*\beta$ equals
$-\frac12[D]^2$. In ordinary cohomology
of $M\times M$, dimensional reasons give
$[A]^2=[B']^2=0$ and, since $n$ is even,
$[A]\smile[B']=[B']\smile[A]$. Consequently
\[
 -\frac12[D]^2=[A]\smile[B']
 \quad\text{in }H^{2n}(M\times M;\R),
\]
so $\langle\pi^*[\beta],[M\times M]\rangle=1$.
The same bar comparison and duality as before
now yield
\begin{equation}\label{eq:even-bound}
 \|A_M\|_\Delta\geq
 \frac{2\|M\|^2}{n^2((n+1)!)^2}.
\end{equation}

\begin{proof}[Proof of Theorem~\ref{thm:intro-estimate}]
The upper bound is Lemma~\ref{lem:upper-norm}.
The lower bound is \eqref{eq:odd-bound}
when $n$ is odd and \eqref{eq:even-bound}
when $n$ is even. Since $\|M\|>0$,
the class $A_M$ is nonzero. Thus
$\mathfrak a_{2n}$ is nonzero as well.
\end{proof}

\begin{remark}
The constants are explicit rather than optimal.
Our argument does not show
$\|A_M\|_\Delta=\|M\times M\|$.
In particular, it does not prove equality
between the cellular seminorm of the nerve
and the simplicial volume of the product.
\end{remark}


\section{Generic metrics and geodesic complexity}
\label{sec:generic}

We establish the general-position statements used above and
give an upper bound for geodesic complexity.
Throughout, $N=2n=\dim(M\times M)$.
All metrics in this section have negative sectional curvature.

We begin by recording the smoothness and variation properties of the distance branches.

Let $g$ be a negatively curved metric on $M$.
The lifted metric on $X=\widetilde M$ is complete,
simply connected and has no conjugate points.
For $(u,v)$ with $u\ne av$, the distance branch
$L_a(u,v)=d_X(u,av)$ is smooth.
Moreover, in a sufficiently small finite-dimensional
smooth family of negatively curved metrics, the branches
$L_a$ vary smoothly in the metric parameters and endpoints:
the unique connecting geodesics vary smoothly by the
geodesic equation and nondegeneracy of the endpoint map.
We use this smoothness only off the diagonal $x=y$,
where different minimizing labels cannot tie.

Two distinct minimizing geodesics from $x$ to $y$
have disjoint interiors. Indeed, if they shared an
interior point, swapping their initial and final
segments would again give a shortest broken path.
A globally minimizing path cannot have a corner,
so uniqueness for the geodesic equation forces
the geodesics to coincide.

\begin{lemma}\label{lem:metric-submersion}
Suppose $a_0,\ldots,a_{k-1}$ represent distinct
minimizing geodesics from $x$ to $y$.
There are smooth variations of the metric,
supported in pairwise disjoint neighborhoods
of interior points of these geodesics, whose
first variations of their lengths form an
invertible diagonal matrix. Consequently,
the universal map
\[
 (g,u,v)\longmapsto
 (L_{a_1}-L_{a_0},\ldots,L_{a_{k-1}}-L_{a_0})
\]
is a submersion at a simultaneous minimum.
\end{lemma}

\begin{proof}
Choose an interior point of each geodesic not
lying on any other minimizing geodesic, and
a nonnegative smooth cutoff $\chi_i$ supported
in a small neighborhood of that point.
Use the symmetric two-tensor variation
$h_i=\chi_i g$. The first variation of the
length of a geodesic with fixed endpoints is
\[
 D_g L_{a_j}(h_i)
 =\frac12\int_{\gamma_j}
    h_i(\dot\gamma_j,\dot\gamma_j)\,ds.
\]
The integral is positive for $i=j$ and
vanishes for $i\ne j$, by disjointness of
the supports. Taking the differences with
the $a_0$ branch gives independent
variations in the target.
\end{proof}

These variations lead to a general-position theorem.

\begin{theorem}\label{thm:genericity}
The set of regular metrics is open and dense
in $\mathcal R_-(M)$.
For a regular metric, the locus of pairs with
exactly $k$ minimizing geodesics is locally a
smooth submanifold of codimension $k-1$.
No such pairs exist when $k>N+1$, and the
locus with $k=N+1$ is finite.
\end{theorem}

\begin{proof}
Fix $g_0\in\mathcal R_-(M)$ and a compact set
$K\subset\widetilde M$ whose deck translates cover
$\widetilde M$. In a sufficiently small $C^\infty$
neighborhood $\mathcal U$ of $g_0$, the lifted metrics
are uniformly bilipschitz to $g_0$, with constant
arbitrarily close to one. Properness of the deck action
and compactness of $K$ therefore give a finite set
$F\subset\Gamma$ such that every minimizing label
at a point of $K\times K$, for every $g\in\mathcal U$,
belongs to $F$. Enlarge $F$, if necessary, to include
all labels whose lengths differ from the minimum by at
most a fixed small allowance. All possible simultaneous
minima are thus drawn from finitely many subsets of $F$.

For a subset $S=\{a_0,\ldots,a_{k-1}\}\subset F$ with
$k\geq2$, consider on the locus where the branches
are smooth the universal difference map
\[
 \mathcal F_S(g,u,v)
 =\bigl(L_{a_1}^g-L_{a_0}^g,\ldots,
        L_{a_{k-1}}^g-L_{a_0}^g\bigr).
\]
At a point where all labels in $S$ are minimizing,
Lemma~\ref{lem:metric-submersion} makes
$\mathcal F_S$ a submersion in the metric parameter
directions. This property is open: after choosing
the finitely many metric variations supplied by the
lemma, their length-variation matrix remains
surjective on a neighborhood of that point.

For the fixed base metric $g_0$, the simultaneous
minimum locus associated with $S$ in $K\times K$ is
compact. Cover it by finitely many such neighborhoods
and collect the metric variations from these
neighborhoods. Their span is a finite-dimensional
space of smooth symmetric two-tensors. Shrink the
metric-parameter ball so that the resulting metrics
remain negatively curved. On a neighborhood of the
simultaneous-minimum locus, the corresponding
finite-parameter map $\mathcal F_S$ is transverse
to zero. On the compact complement of that
neighborhood, simultaneous minimization for $S$
is excluded by a strictly positive length gap;
after a further shrinking of the parameter ball,
it remains excluded. Thus finite-dimensional
parametric transversality produces arbitrarily
small perturbations for which the fixed-label
difference map is transverse to zero at every
simultaneous minimum for $S$.

Repeat the construction for the finitely many
subsets $S\subset F$ and take the span of all
the resulting metric variations. Transversality
for one configuration is stable under adjoining
the other parameters, because the earlier
surjective parameter directions remain available.
The parametric transversality theorem, applied
to all configurations simultaneously, gives
metrics arbitrarily close to $g_0$ satisfying
the regularity condition. This proves density.

For openness, suppose $g_j\to g$ and that every
$g_j$ is nonregular, whereas $g$ is regular.
After putting the corresponding lifts of endpoint
pairs in $K\times K$, the finite-label reduction
allows us to pass to a subsequence with a fixed
set $S$ of at least two minimizing labels witnessing
failure of full rank. Compactness gives a limit
$(u,v)$ at which all labels in $S$ still minimize
for $g$. The non-full-rank condition is closed
under convergence of the differentials, contradicting
regularity of $g$. Hence regular metrics form an
open subset.

At a regular pair with exactly $k$ minimizing
labels, the implicit function theorem cuts out
the equal-length locus with codimension $k-1$.
All other labels have a positive length gap in a
small neighborhood, so this is also the local
exact-multiplicity stratum. The inequality
$k-1\leq N$ follows from full rank. The stratum
of multiplicity $N+1$ has dimension zero.
It is closed, since the existence of at least
$N+1$ minimizing labels is a closed condition
under the finite-label reduction; compactness
of $M\times M$ then makes it finite.
\end{proof}

We next establish a selection result for minimizing branches on the multiplicity strata.

We recall that $\GC(M,g)$ is the minimum
number of locally compact pieces partitioning
$B=M\times M$, each with a continuous choice
of minimizing geodesic, using the unreduced
convention of \cite{RecioMitter2021}.
The following selection argument is useful
more generally for a finite-valued covering
with clean multiplicity strata.

\begin{lemma}\label{lem:selection}
Let $B$ be a second-countable smooth $N$-manifold
(without boundary), let $q:E\to B$ be a smooth
covering, and let $Z\subset E$ be closed. Assume
$q|_Z$ is proper and surjective, with nonempty
finite fibers. For $k\geq1$ put
\[
 B_k=\{b\in B:\#(Z\cap q^{-1}(b))=k\}.
\]
Assume every nonempty $B_k$ is a smooth
embedded submanifold of $B$ (not necessarily
closed), of dimension at most $N-k+1$.
Then $B$ admits a partition into at most
$N+1$ pairwise disjoint locally closed subsets
on each of which $q|_Z$ admits a continuous
section. In particular, the conclusion does
not assume that the finite-valued map
$q|_Z$ is a covering.
\end{lemma}

\begin{proof}
Because $q|_Z$ is proper and $Z$ is closed,
the fiber cardinality $k(b)$ is upper
semicontinuous. Indeed, given $b\in B$, choose
an evenly covered neighborhood and disjoint
local sheets through the finitely many points
of $Z_b$. Properness excludes any other
$Z$-points over a sufficiently small
neighborhood of $b$. Thus nearby fibers
are subsets of these finitely many sheets.
Consequently $B_k$ is locally closed;
the smooth embedded structure assumed in
the statement is compatible with the
subspace topology. Moreover $k\leq N+1$
whenever $B_k$ is nonempty.

Write $d_k=\dim B_k\leq N-k+1$. We first
choose $N+1$ smooth functions
$f_0,\ldots,f_N:E\to\R$ with a simultaneous
general-position property on all the strata.

Fix $k$ and a point of $B_k$. Since $q$ is a
covering and $Z\to B$ is proper with finite
fibers, after shrinking a neighborhood of the
point we can enumerate the $k$ active sheets
by smooth local sections of $q$. Restricting
these sections to $B_k$, we obtain smooth
maps $z_1,\ldots,z_k$ with pairwise disjoint
images in $E$. For a tuple
$f=(f_0,\ldots,f_N)$, consider the evaluation
map
\[
 F_f:B_k\longrightarrow (\R^k)^{N+1},
 \qquad
 F_f(b)=\bigl((f_j(z_i(b)))_{i=1}^{k}
                  \bigr)_{j=0}^{N}.
\]
For each collection
$\mathcal P=(\mathcal P_0,\ldots,\mathcal P_N)$
of partitions of $\{1,\ldots,k\}$, let
$D_{\mathcal P}\subset(\R^k)^{N+1}$ be
the linear subspace imposing equality of
coordinates inside each block of
$\mathcal P_j$, for each $j$. Its codimension
is
\[
 \operatorname{codim}D_{\mathcal P}
 =\sum_{j=0}^N(k-|\mathcal P_j|).
\]
The parametrized evaluation map
$(f,b)\mapsto F_f(b)$ is a submersion in the
function variables: values at the finitely
many distinct points $z_i(b)$ can be
perturbed independently with disjointly
supported smooth bump functions on $E$.
Parametric transversality therefore gives
a residual set of tuples $f$ for which
$F_f$ is transverse to every
$D_{\mathcal P}$ on the chosen chart.

There are only finitely many partitions for
each $k$, and countably many sheets and
charts suffice to cover the second-countable
strata $B_k$. Using compact exhaustions
and the Baire property of the space of
smooth functions on $E$, choose one tuple
$f$ transverse to all these diagonals
simultaneously. When
$\operatorname{codim}D_{\mathcal P}>d_k$,
transversality implies that
$F_f^{-1}(D_{\mathcal P})$ is empty.
Consequently, at every $b\in B_k$, if
$\mathcal P_j(b)$ is the partition of the
active $k$ sheets according to equality of
their $f_j$-values, then
\begin{equation}\label{eq:selection-inequality}
 \sum_{j=0}^N
 \bigl(k-|\mathcal P_j(b)|\bigr)
 \leq d_k\leq N-k+1.
\end{equation}

Define
\[
 c(b)=k(b)-1+
 \sum_{j=0}^N
 \bigl(k(b)-|\mathcal P_j(b)|\bigr).
\]
Thus $0\leq c(b)\leq N$.
The vectors
\[
 (f_0(z),\ldots,f_N(z)),\qquad z\in Z_b,
\]
are pairwise distinct. Otherwise two
active sheets would have equal values for
all $N+1$ functions, forcing each term
$k-|\mathcal P_j|$ to be at least one,
contrary to \eqref{eq:selection-inequality}.
Hence there is a unique lexicographically
smallest vector, determining a selected
sheet over each $b$.

We verify that the selection is continuous
on the sets $A_j=\{b:c(b)=j\}$.
Fix $b$ and enumerate all $k(b)$ active
sheets in a covering neighborhood $U$.
Properness of $Z\to B$ and closedness of
$Z\subset E$ allow $U$ to be chosen so
that every active sheet at $b'\in U$
is among these enumerated sheets.
Furthermore, pairs with unequal
$f_j$-values at $b$ remain unequal after
shrinking $U$. Therefore, as $b'\to b$,
the number of active sheets cannot
increase, and each equality partition can
only split or lose elements. More precisely,
for every $j$ the nonnegative integer
$k(b')-|\mathcal P_j(b')|$ cannot exceed
its value at $b$, and $k(b')\leq k(b)$.
Hence $c$ is upper semicontinuous.

Since $c$ is integer-valued with range
$\{0,\ldots,N\}$, each
\[
 A_j=\{c\geq j\}\cap\{c<j+1\}
\]
is locally closed and thus locally compact.
Fix $b$ and shrink $U$ as above. For
$b'\in A_{c(b)}\cap U$, all nonnegative
contributions
\[
 k(b')-1,\qquad
 k(b')-|\mathcal P_j(b')|
 \quad(0\leq j\leq N)
\]
are bounded above by the corresponding
ones at $b$. Their sums are equal, so
each contribution agrees. Hence all
$k(b)$ active sheets persist and the
equality partitions have not split.
Every pair of sheet values that was
unequal at $b$ remains unequal with the
same sign in a sufficiently small
neighborhood. Thus the lexicographic
order between any two distinct evaluation
vectors is constant along
$A_{c(b)}\cap U$. Its unique minimum
therefore belongs to the same covering
sheet as at $b$, proving continuity
of the selected section on $A_{c(b)}$.
The disjoint locally closed sets
$A_0,\ldots,A_N$ therefore give the
required partition and continuous
sections.
\end{proof}

\begin{theorem}\label{thm:generic-GC}
For every regular negatively curved metric
on $M$ one has
\[
 \GC(M,g)\leq 2n+1.
\]
If $M$ is closed and oriented, then
\[
 \GC(M,g)=\TC(M)=2n+1
\]
and
\[
 \GC(M,g)
 =1+\max_{x\ne y}
   \operatorname{rank}\MH_{2,d_M(x,y)}(M;x,y).
\]
\end{theorem}

\begin{proof}
Let $E\to B$ be the covering whose fiber
over $(x,y)$ consists of fixed-endpoint
homotopy classes of paths, equivalently
the labels of lifts. Let $Z$ be the subset
corresponding to minimizing geodesics.
The map $Z\to B$ is closed and proper with
finite fibers, by compactness of $M$ and
properness of the universal covering action.
The exact multiplicity strata are smooth
of dimension $N-k+1$, by
Theorem~\ref{thm:genericity}.
Apply Lemma~\ref{lem:selection}. A continuous
selection of a minimizing label gives a
continuous minimizing geodesic, because
the corresponding lifted geodesic depends
smoothly on its endpoints. Thus
$\GC(M,g)\leq N+1$.

The inequality $\TC(M)\leq\GC(M,g)$ is
standard \cite{RecioMitter2021}. The
known calculation of topological complexity
for torsion-free hyperbolic groups
\cite{Dranishnikov2020}, translated into
our unreduced convention, gives
$\TC(M)=N+1$. Hence equality holds.

For the final assertion, we do not use the
quantitative counting theorem. Horja's
multiplicity theorem \cite{Horja1997}
shows that a pair realizing a local maximum
of the distance function on a closed
negatively curved $n$-manifold is joined
by at least $2n+1=N+1$ minimizing
geodesics. Such a pair exists by compactness.
For a regular metric, Theorem~\ref{thm:genericity}
gives the reverse multiplicity bound $N+1$.
Hence the maximal number of minimizing
geodesics is exactly $N+1$. By
Proposition~\ref{prop:MH-minimal}, the
largest minimal-length degree-two magnitude
homology rank is $N$.

In particular, the geodesic-complexity
calculation and this rank identity are
independent of the quantitative seminorm
and localization estimates in
Sections~\ref{sec:seminorm} and
\ref{sec:localization}.
\end{proof}


\section{Localization and the counting theorem}
\label{sec:localization}

In this section $M$ is closed, oriented and negatively curved,
$N=2n$, and $g$ is a regular negatively curved metric.
We prove the localization formula announced in
Section 1 and the quantitative counting theorem.

We first calculate the local degree at a point of maximal multiplicity.

Let $\mathcal N_g$ be the minimizing-geodesic nerve
constructed in Section~\ref{sec:nerve}, and let
$Q_g=|\mathcal N_g|/G$.
By Theorem~\ref{thm:genericity}, $\dim\mathcal N_g\leq N$.
The equivariant map $\Phi_g$ from
Proposition~\ref{prop:Phi} descends to
$\bar\Phi_g:B\to Q_g$.

For each $b=(x,y)\in C_g$, choose lifts
$(u,v)$ and an ordered minimizing-label set
\[
 S_b=\{a_0,\ldots,a_N\}.
\]
The associated top-dimensional simplex of
$\mathcal N_g$ is denoted $\sigma_b$.
Write
\[
 F_b=(L_{a_1}-L_{a_0},\ldots,L_{a_N}-L_{a_0}).
\]
The regularity condition says that
$dF_b:T_bB\to\R^N$ is invertible.
Define
\begin{equation}\label{eq:local-sign}
 \epsilon_b=\operatorname{sgn}\det dF_b\in\{+1,-1\}.
\end{equation}

\begin{lemma}\label{lem:local-degree}
The local degree of $\bar\Phi_g$ at $b$,
with target orientation determined by the
ordered simplex $[a_0,\ldots,a_N]$, equals
$\epsilon_b$.
\end{lemma}

\begin{proof}
Choose a small neighborhood of $b$ and compatible
lifts so that the only labels with
$L_a<m+\delta$ are the labels in $S_b$.
This is possible after shrinking the neighborhood
and the uniform allowance $\delta$, without
changing the nerve.

On this neighborhood the barycentric weights are
\[
 \lambda_i=\frac{m+\delta-L_{a_i}}{
       \sum_{j=0}^{N}(m+\delta-L_{a_j})},
 \qquad 0\leq i\leq N.
\]
All weights are strictly positive near $b$.
Their differences satisfy the exact identity
\begin{equation}\label{eq:barycentric-difference}
 \lambda_i-\lambda_0
 =-\frac{L_{a_i}-L_{a_0}}{W},
 \qquad 1\leq i\leq N.
\end{equation}
The barycenter is characterized by the vanishing
of the $N$ differences in
\eqref{eq:barycentric-difference}, so its inverse
image near $b$ is precisely $\{b\}$.
At $b$, $W=(N+1)\delta>0$. Although
$m=\min_a L_a$, and hence $W$, need not be
smooth at $b$, the continuous positive factor
$1/W$ does not affect the local Brouwer degree
of the map in \eqref{eq:barycentric-difference}.
Indeed, in local coordinates supplied by the
inverse function theorem for $F_b$, the
difference-coordinate map is
$t\mapsto-t/W(F_b^{-1}(t))$.
On a sufficiently small ball about the origin,
this is homotopic through maps with no boundary
zeros to $t\mapsto-t/W(b)$, by interpolating the
positive denominator with the constant $W(b)$.

The orientation defined by
$(\lambda_1-\lambda_0,\ldots,\lambda_N-\lambda_0)$
agrees with the standard simplex orientation:
the transition matrix from
$(\lambda_1,\ldots,\lambda_N)$ has
positive determinant $N+1$.
Therefore the local degree equals
$(-1)^N\operatorname{sgn}\det dF_b
=\operatorname{sgn}\det dF_b$, since $N=2n$.
No smoothness of $\Phi_g$ across the
minimum-length strata is required.
\end{proof}

\begin{lemma}\label{lem:localized-cycle}
The cellular cycle
\begin{equation}\label{eq:nerve-cycle}
 Z_g=\sum_{b\in C_g}\epsilon_b[\sigma_b]
 \in C_N(Q_g;\Z)
\end{equation}
represents $(\bar\Phi_g)_*[B]$ in $H_N(Q_g;\Z)$.
If multiple points yield the same orbit of
oriented simplices, their coefficients are added.
\end{lemma}

\begin{proof}
By Lemma~\ref{lem:no-inversions}, each open top
cell of $Q_g$ has an ordinary characteristic
simplex without vertex permutations. For each open top cell choose its barycenter.
Its inverse images are isolated points in $C_g$,
and Lemma~\ref{lem:local-degree} computes their
Brouwer local degrees. The cellular coefficient
of the pushforward of the fundamental class
is the sum of these degrees, by the relative
homology definition of degree. This reasoning
uses continuous maps and does not require
the barycenter to be a smooth regular value.
The resulting chain is \eqref{eq:nerve-cycle}.
It is a cycle because $B$ has no boundary.
\end{proof}

We now assemble the local contributions into a global formula.

\begin{theorem}\label{thm:localization}
Let $M$ be connected, closed, oriented and negatively curved,
and let $g$ be regular. Then
\begin{equation}\label{eq:localization}
 \mathfrak a_N\frown[M\times M]
 =\sum_{b\in C_g}\epsilon_b
 \big[(a_1-a_0)\wedge\cdots\wedge(a_N-a_0)\big]
 \quad\text{in }(\bigwedge^N I)_G.
\end{equation}
The right-hand side does not depend on the
chosen lifts or vertex orderings, provided
$\epsilon_b$ is defined by \eqref{eq:local-sign}.
\end{theorem}

\begin{proof}
Compose $\bar\Phi_g:B\to Q_g$ with the
simplicial inclusion $Q_g\to Q$, where
$Q=|Y|/G$ is the full simplex orbit space.
Proposition~\ref{prop:comparison} gives
\[
 \iota_N(\bar\Phi_g)_*[B]
 =\mathfrak a_N\frown[B].
\]
By Lemma~\ref{lem:localized-cycle}, the class
on the left is represented by $Z_g$.
Lemma~\ref{lem:exterior-quotient} sends
the simplex $[a_0,\ldots,a_N]$ to
\[
 (a_1-a_0)\wedge\cdots\wedge(a_N-a_0),
\]
which proves the formula.

Changing lifts applies the same $G$-element
to all labels, hence does not alter their
coinvariant class. Permuting the labels
changes the orientation of the simplex and
the sign of the determinant by the same
sign. The product appearing in
\eqref{eq:localization} is therefore unchanged.
\end{proof}

\begin{proposition}\label{prop:endpoint-reversal}
Let $J:\bigl(\bigwedge^N I\bigr)_G\to
\bigl(\bigwedge^N I\bigr)_G$ be the involution
induced by the additive map
\[
 \Z[\Gamma]\longrightarrow\Z[\Gamma],
 \qquad a\longmapsto a^{-1}.
\]
For $b=(x,y)\in C_g$, write
\[
 \tau_b=\epsilon_b
 \big[(a_1-a_0)\wedge\cdots\wedge(a_N-a_0)\big].
\]
Then the reversed pair $b^{\mathrm{op}}=(y,x)$
also belongs to $C_g$, and
\begin{equation}\label{eq:endpoint-reversal}
 \tau_{b^{\mathrm{op}}}=(-1)^n J(\tau_b).
\end{equation}
In particular $\#C_g$ is even and
\[
 J(A_M)=(-1)^n A_M.
\]
More precisely, $A_M$ lies in the image of
$1+(-1)^n J$ on the integral coinvariant module.
\end{proposition}

\begin{proof}
The inversion map sends the action
$(g,h)\cdot a=gah^{-1}$ to the action with
the two group coordinates interchanged:
\[
 (gah^{-1})^{-1}=h a^{-1}g^{-1}.
\]
It therefore induces the stated involution
on coinvariants and their exterior powers.

Choose lifts $u,v$ of $x,y$, and label the
minimizing geodesics by
$a_0,\ldots,a_N$. For the reversed endpoints
choose lifts $v,u$ and labels
$a_0^{-1},\ldots,a_N^{-1}$. The equality
\[
 d_X(v,a_i^{-1}u)=d_X(u,a_iv)
\]
shows that the reversed length-difference
map is the original length-difference map
composed with the endpoint swap
$T(u,v)=(v,u)$. On the $N=2n$ dimensional
product, $T$ has orientation degree
$(-1)^{n^2}=(-1)^n$. Hence the two
determinant signs satisfy
\[
 \epsilon_{b^{\mathrm{op}}}
 =(-1)^n\epsilon_b.
\]
Combining this with inversion of the
exterior generator gives
\eqref{eq:endpoint-reversal}.

The diagonal consists only of pairs joined
by one minimizing geodesic, so reversal acts
freely on $C_g$. The cardinality is even.
Summing \eqref{eq:endpoint-reversal}
over all points and using the localization
formula proves the eigenspace identity.
Selecting one point from each reversal
orbit also expresses $A_M$ as
$(1+(-1)^nJ)$ applied to an integral element.
\end{proof}

The local exterior terms admit a description in magnitude homology.

At each $b=(x,y)\in C_g$ the label set has
$N+1$ elements. Proposition
\ref{prop:MH-minimal} gives
\[
 \MH_{2,d_M(x,y)}(M;x,y)\cong I(S_b)
\]
with $\operatorname{rank}I(S_b)=N$.
The inclusion $I(S_b)\hookrightarrow I$
takes the generator
\[
 (a_1-a_0)\wedge\cdots\wedge(a_N-a_0)
 \in\bigwedge^N I(S_b)
\]
to the local exterior term in
\eqref{eq:localization}.
This interpretation is canonical up to the
orientation sign already accounted for
by $\epsilon_b$. Thus each summand is a
local index valued in the top exterior
power of second magnitude homology at
the minimal endpoint distance.

Combining localization with the seminorm estimate proves the counting theorem.

\begin{proof}[Proof of Theorem~\ref{thm:intro-count}]
Theorem~\ref{thm:genericity} proves that
regular metrics form an open dense set
and that $C_g$ is finite for every such metric.
By \eqref{eq:localization}, the real class
$A_M$ has a representation as a sum of
$\#C_g$ exterior generators, each with
coefficient $+1$ or $-1$. Definition
\eqref{eq:delta-norm} therefore yields
\[
 \|A_M\|_\Delta\leq\#C_g.
\]
The lower estimate from
Theorem~\ref{thm:intro-estimate} gives
\[
 \#C_g\geq\|A_M\|_\Delta
 \geq c_n\|M\|^2.
\]
The claim follows.
\end{proof}

\begin{corollary}\label{cor:top-multiplicity}
For every regular negatively curved metric on an
oriented closed negatively curved $n$-manifold,
there exists an ordered pair of points joined
by exactly $2n+1$ minimizing geodesics.
\end{corollary}

\begin{proof}
The simplicial volume of a closed negatively
curved manifold is positive, so the lower
bound in Theorem~\ref{thm:intro-count}
forces $\#C_g>0$.
\end{proof}

\begin{example}\label{ex:surface-count}
For a closed oriented surface $\Sigma_h$ of genus
$h\geq2$, one has $\|\Sigma_h\|=4(h-1)$.
Here $n=2$ and $c_2=1/72$. Therefore
every regular negatively curved metric satisfies
\[
 \#\{(x,y):m_g(x,y)=5\}
 \geq\frac{2}{9}(h-1)^2.
\]
The count concerns ordered pairs.
\end{example}

\begin{remark}
The estimate does not state that every negatively
curved metric is regular. At a nongeneric metric,
several top-dimensional local contributions may
coalesce at a point of higher multiplicity.
A general local-index formula for such degeneracies
requires additional analysis and is not needed
for the counting theorem proved here.
\end{remark}


\section{Further remarks and questions}
\label{sec:remarks}

We conclude with several distinctions and open questions
arising from the results established above. The emphasis
is on identifying what the present argument does and does
not imply.

\subsection{The exterior class and ordinary topological complexity}

The ordinary tensor-valued Costa--Farber canonical class
$\mathfrak v^{2n}$ is an obstruction to motion planning
\cite{CostaFarber2010}. The exterior class
$\mathfrak a_{2n}$ is obtained from it by a quotient of
coefficient modules. Theorem~\ref{thm:intro-estimate}
shows that this exterior image is nonzero for closed
oriented negatively curved manifolds. This is useful
because the exterior coefficient module is a quotient
of the oriented simplicial chains of the full simplex.

However, the exterior construction does not give a
new value of $\TC(M)$. The equality $\TC(M)=2n+1$
was already known \cite{Dranishnikov2020}.
The additional content here is a quantitative
connection with multiplicity and a localization
formula whose terms are detected by
minimal-length magnitude homology.

\subsection{Comparison with cut-locus methods}

Geodesic complexity is sensitive to the structure
of the total cut locus. Existing work provides
upper bounds from decompositions of cut loci
\cite{RecioMitter2021,MescherStegemeyer2023}.
Our regularity condition instead controls all
simultaneously minimizing length branches,
including the highest multiplicity stratum.
The selection construction in
Section~\ref{sec:generic} uses this transverse
stratification to obtain an upper bound.

There is no general inequality identifying
geodesic complexity with the maximum number
of minimizing geodesics. Their equality in
Theorem~\ref{thm:generic-GC} is a consequence
of the specific geometric and cohomological
hypotheses, not a formal property of
geodesic complexity.

\subsection{The strength of the seminorm}

The seminorm $\|\cdot\|_\Delta$ is defined
by a concrete generating set in the
exterior augmentation coinvariants.
Our estimates show
\[
 c_n\|M\|^2
 \leq\|A_M\|_\Delta
 \leq\|M\times M\|.
\]
The inequalities need not be equalities.
Determining the exact value of
$\|A_M\|_\Delta$ even for closed hyperbolic
surfaces appears to require a more precise
description of bounded alternating
cochains on $\Gamma$.

The constants $c_n$ arise from the
norm-controlled lifting procedure.
They are not expected to be sharp in
general. An improvement would follow
from a smaller lifting norm for classes
built from the orientation cocycle,
without requiring a general isometric
isomorphism in bounded cohomology.

\subsection{Nongeneric metrics}

For a regular metric the highest
multiplicity stratum is finite and each
local contribution to the localization
formula has absolute coefficient one.
For a nongeneric metric one may have
isolated points with more than $2n+1$
minimizing geodesics, or positive-dimensional
sets of degenerate configurations.

The nerve construction remains available
for any negatively curved metric, but
the localization as a finite signed sum
of top-stratum points uses regularity.
A possible extension is to use relative
homology of the local simplex and its
lower-dimensional skeleton to define
local indices at isolated degenerate
configurations. Such an extension is
not proved here.

\subsection{Magnitude homology beyond degree two}

The identification
\[
 \MH_{2,d_M(x,y)}(M;x,y)\cong I(S_g(u,v))
\]
uses two special features: the length
is exactly the endpoint distance,
and the homological degree is two.
Higher magnitude homology groups
incorporate more complicated chains
and need not be determined only by
the cardinality of the minimizing
geodesic set. Their relationship
with higher-dimensional strata of
the minimizing-geodesic nerve is
a separate question.

The comparison with interval-poset
models in \cite{KanetaYoshinaga2021}
is a natural starting point, but
the existence of a global sheaf or
derived-category refinement is not
asserted in this paper.


\bibliographystyle{plain}
\bibliography{references}

\end{document}